\section{Trade: The Distance Tariff}
\label{sec:trade}

Containerisation is the archetypal equaliser: it slashed the cost of moving
goods between big ports. But the container economy rewards scale. Shipping
lines deploy their largest vessels on the busiest routes and serve small
markets with fewer, smaller and less frequent calls, often through
transhipment. For an island importing most of its food, fuel and building
materials, that means higher freight rates, longer and less predictable
lead times, and larger inventories---a tariff levied by geography rather than
by any government.

UNCTAD's Liner Shipping Connectivity Index (LSCI) measures how well a
country is plugged into the global liner network, combining the number of
scheduled calls, the capacity deployed, the number of services and companies,
the size of the largest ships and the number of direct connections \parencite{unctad2026lsci}.
\Cref{fig:lsci} tracks it from \lsciFirstYear{} to \lsciLastYear.

\begin{figure}[t]
\centering
\begin{tikzpicture}
\begin{axis}[edge, height=0.55\linewidth,
    ymode=log, ymin=10, ymax=800,
    ytick={10,20,50,100,200,500}, yticklabels={10,20,50,100,200,500},
    xmin=2006, xmax=2025, xtick={2006,2010,2015,2020,2025},
    x tick label style={/pgf/number format/1000 sep={}},
    ylabel={Liner Shipping Connectivity Index (log scale)},
    legend pos=north east, legend columns=1]
  \addplot[bandTeal, forget plot, name path=hi] table[x=year,y=sids_p75] {\datadir/lsci_groups.csv};
  \addplot[bandTeal, forget plot, name path=lo] table[x=year,y=sids_p25] {\datadir/lsci_groups.csv};
  \addplot[teal, fill opacity=0.12, fill=teal, draw=none, forget plot] fill between[of=hi and lo];
  \addplot[lineEmber] table[x=year,y=hubs] {\datadir/lsci_groups.csv};
  \addlegendentry{2006's ten best-connected hubs (median)}
  \addplot[lineSlate] table[x=year,y=all] {\datadir/lsci_groups.csv};
  \addlegendentry{All countries (median)}
  \addplot[lineTeal] table[x=year,y=sids] {\datadir/lsci_groups.csv};
  \addlegendentry{Small island states (median, band = middle half)}
\end{axis}
\end{tikzpicture}
\caption{\takeaway{Shipping connectivity at the edge has flat-lined.}
  The median small island state scored \sidsMedianLSCIfirst{} in
  \lsciFirstYear{} and \sidsMedianLSCI{} in \lsciLastYear{}, while the median
  country rose by \allLSCIGrowthPct\%. The old hubs lost ground to new ones, so the
  hub-to-island ratio narrowed from \hubSidsRatioFirst{} to \hubSidsRatio{} times.
  Annual averages of quarterly data; index average across countries in Q1 2023 = 100.
  \source{UNCTAD Liner Shipping Connectivity Index; UN M49 SIDS list.}}
\label{fig:lsci}
\end{figure}


The data tell a more nuanced story than the thesis predicts. The headline is
\emph{stagnation at the edge}: the median small island state scored
\sidsMedianLSCIfirst{} in \lsciFirstYear{} and \sidsMedianLSCI{} in
\lsciLastYear{}, essentially unchanged across two decades in which the
median country's connectivity grew by \allLSCIGrowthPct\% and global container
trade expanded. Relative to the median country, island states fell
slightly further behind: the ratio went from \allSidsRatioFirst{} to
\allSidsRatio. Of \nSIDSlsci{} island states with data at both ends,
\nSIDSlsciFell{} saw their index fall. Tuvalu scores \lsciTuvalu; Tonga
\lsciTonga, ranking \lsciTongaRank{} of \nLsci{} economies.

But the gap with the very top did \emph{not} widen. The ten hubs that were
best connected in \lsciFirstYear{} (\hubList) saw their median index fall
from \hubMedianLSCIfirst{} to \hubMedianLSCI{}, as connectivity shifted
towards China and other emerging hubs. The hub-to-island ratio therefore
narrowed, from \hubSidsRatioFirst{} to \hubSidsRatio{} times. A claim that
islands are being left ever further behind the shipping superpowers is not
supported; the more defensible claim is that the global network grew denser
almost everywhere \emph{except} at its smallest, most remote nodes.

\begin{counterpoint}[Islands can climb]
Mauritius raised its LSCI from \lsciMauritiusFirst{} to \lsciMauritius{} over
the period, partly by developing Port Louis as a transhipment centre, and
Singapore built connectivity ranked \lsciSingaporeRank{} of \nLsci{} on a small island
(\cref{sec:escape}). Shipping connectivity is shaped by port investment and
policy, not just by geography.
\end{counterpoint}

The distance tariff also shows up in what islands can export. Goods that are
heavy, perishable or time-sensitive are penalised most, steering remote
economies towards tourism, fish, remittances and rents---activities that
travel by aeroplane, satellite or bank transfer rather than by box. That is
an adaptation, not a failure, but it concentrates risk: when a pandemic grounds
flights or a cyclone closes the port, there are few alternatives.
